% GENERATED FILE. Edit canonical lessons, then run npm run generate:books.
% canonical-source-hash: fnv1a64:5db330be71cc4e05
% canonical-lessons: ES-C741-cansancio, ES-C741-agotado, ES-C741-dormirse, ES-C741-insomnio, ES-C741-siesta

\chapter{Cansancio y descanso}
\label{ch:es-cansancio-y-descanso}
\hlchaptermodality{741}

\begin{chapteropening}
\textbf{By the end of this chapter:} \emph{I can say how tired I am and how I sleep.}

Complete the last lesson of chapter 741: I can say how tired I am and how I sleep.
\end{chapteropening}

\section[el cansancio]{el cansancio — tiredness}

\label{lesson:ES-C741-cansancio}



\begin{quote}
Before the new one: say the Spanish for to faint, then the Spanish for the allergy.
\end{quote}

\subsection*{You'll want to know: el cansancio}
\textbf{el cansancio} — \textquotedblleft{}tiredness\textquotedblright{}.

\begin{quote}
Siento mucho cansancio.
\end{quote}

\begin{quote}
\emph{I feel very tired.}
\end{quote}

\begin{cousinweb}
\emph{cansancio} is built on \emph{cansar} \textquotedblleft{}to tire\textquotedblright{}, which comes from Latin \emph{campsāre} \textquotedblleft{}to bend, to sail round a headland\textquotedblright{} (EtymDB; Lewis \& Short). How the sense moved to tiring is not settled.
\end{cousinweb}

\begin{grammarlens}[title={What you've built}]
One more word for telling, explaining and describing at B1.
\end{grammarlens}

\subsection*{Your turn}
\begin{itemize}
\raggedright
  \item \emph{Say it:} \emph{el cansancio}
  \item \emph{Say it:} \emph{el cansancio}, once more
  \item \emph{Say it:} a sentence of your own with \emph{el cansancio}
  \item \emph{Recall:} say the Spanish for to faint, then the Spanish for the allergy, then say \emph{el cansancio} again
\end{itemize}

\begin{tcolorbox}[breakable,colback=teal!4,colframe=teal!35!black,title={Before you move on}]
What does el cansancio mean? (\textquotedblleft{}Tiredness\textquotedblright{}.) Say it once more.
\end{tcolorbox}
\section[agotado]{agotado — exhausted; sold out}

\label{lesson:ES-C741-agotado}



\begin{quote}
Before the new one: say the Spanish for the allergy, then the Spanish for tiredness.
\end{quote}

\subsection*{You'll want to know: agotado}
\textbf{agotado} — \textquotedblleft{}exhausted; sold out\textquotedblright{}.

\begin{quote}
Después del viaje estoy agotado.
\end{quote}

\begin{quote}
\emph{After the trip I am exhausted.}
\end{quote}

On a ticket office: \emph{agotado} — sold out.

\begin{cousinweb}
\emph{agotado} is the participle of \emph{agotar}, usually explained as built on \emph{gota} \textquotedblleft{}drop\textquotedblright{}, from Latin \emph{gutta}: emptied down to the last drop.
\end{cousinweb}

\begin{grammarlens}[title={What you've built}]
One more word for telling, explaining and describing at B1.
\end{grammarlens}

\subsection*{Your turn}
\begin{itemize}
\raggedright
  \item \emph{Say it:} \emph{agotado}
  \item \emph{Say it:} \emph{agotado}, once more
  \item \emph{Say it:} a sentence of your own with \emph{agotado}
  \item \emph{Recall:} say the Spanish for the allergy, then the Spanish for tiredness, then say \emph{agotado} again
\end{itemize}

\begin{tcolorbox}[breakable,colback=teal!4,colframe=teal!35!black,title={Before you move on}]
What does agotado mean? (\textquotedblleft{}Exhausted; sold out\textquotedblright{}.) Say it once more.
\end{tcolorbox}
\section[dormirse]{dormirse — to fall asleep}

\label{lesson:ES-C741-dormirse}



\begin{quote}
Before the new one: say the Spanish for tiredness, then the Spanish for exhausted.
\end{quote}

\subsection*{You'll want to know: dormirse}
\textbf{dormirse} — \textquotedblleft{}to fall asleep\textquotedblright{}.

\begin{quote}
Me dormí viendo la tele.
\end{quote}

\begin{quote}
\emph{I fell asleep watching TV.}
\end{quote}

\emph{Dormirse} also means to oversleep: \emph{me dormí y llegué tarde}.

\begin{cousinweb}
\emph{dormirse} is the reflexive of \emph{dormir}, which is usually traced to Latin \emph{dormīre} \textquotedblleft{}to sleep\textquotedblright{}.
\end{cousinweb}

\begin{grammarlens}[title={What you've built}]
One more word for telling, explaining and describing at B1.
\end{grammarlens}

\subsection*{Your turn}
\begin{itemize}
\raggedright
  \item \emph{Say it:} \emph{dormirse}
  \item \emph{Say it:} \emph{dormirse}, once more
  \item \emph{Say it:} a sentence of your own with \emph{dormirse}
  \item \emph{Recall:} say the Spanish for tiredness, then the Spanish for exhausted, then say \emph{dormirse} again
\end{itemize}

\begin{tcolorbox}[breakable,colback=teal!4,colframe=teal!35!black,title={Before you move on}]
What does dormirse mean? (\textquotedblleft{}To fall asleep\textquotedblright{}.) Say it once more.
\end{tcolorbox}
\section[el insomnio]{el insomnio — insomnia}

\label{lesson:ES-C741-insomnio}



\begin{quote}
Before the new one: say the Spanish for exhausted, then the Spanish for to fall asleep.
\end{quote}

\subsection*{You'll want to know: el insomnio}
\textbf{el insomnio} — \textquotedblleft{}insomnia\textquotedblright{}.

\begin{quote}
Sufro de insomnio.
\end{quote}

\begin{quote}
\emph{I suffer from insomnia.}
\end{quote}

\begin{cousinweb}
\emph{insomnio} is usually traced to Latin \emph{insomnia} \textquotedblleft{}want of sleep\textquotedblright{} (Lewis \& Short), built on \emph{in-} \textquotedblleft{}not\textquotedblright{} and \emph{somnus} \textquotedblleft{}sleep\textquotedblright{}; its ending follows Latin \emph{insomnium}, which in Latin meant a dream.
\end{cousinweb}

\begin{grammarlens}[title={What you've built}]
One more word for telling, explaining and describing at B1.
\end{grammarlens}

\subsection*{Your turn}
\begin{itemize}
\raggedright
  \item \emph{Say it:} \emph{el insomnio}
  \item \emph{Say it:} \emph{el insomnio}, once more
  \item \emph{Say it:} a sentence of your own with \emph{el insomnio}
  \item \emph{Recall:} say the Spanish for exhausted, then the Spanish for to fall asleep, then say \emph{el insomnio} again
\end{itemize}

\begin{tcolorbox}[breakable,colback=teal!4,colframe=teal!35!black,title={Before you move on}]
What does el insomnio mean? (\textquotedblleft{}Insomnia\textquotedblright{}.) Say it once more.
\end{tcolorbox}
\section[la siesta]{la siesta — the nap, the siesta}

\label{lesson:ES-C741-siesta}



\begin{quote}
Before the new one: say the Spanish for to fall asleep, then the Spanish for insomnia.
\end{quote}

\subsection*{You'll want to know: la siesta}
\textbf{la siesta} — \textquotedblleft{}the nap, the siesta\textquotedblright{}.

\begin{quote}
Después de comer echo la siesta.
\end{quote}

\begin{quote}
\emph{After lunch I have a nap.}
\end{quote}

\emph{Echar la siesta} (Spain) or \emph{dormir la siesta}.

\begin{cousinweb}
\emph{siesta} is inherited from Latin \emph{sexta}, short for \emph{sexta hōra} \textquotedblleft{}the sixth hour\textquotedblright{} (EtymDB): counted from dawn, the sixth hour is midday.
\end{cousinweb}

\begin{grammarlens}[title={What you've built}]
One more word for telling, explaining and describing at B1.
\end{grammarlens}

\subsection*{Your turn}
\begin{itemize}
\raggedright
  \item \emph{Say it:} \emph{la siesta}
  \item \emph{Say it:} \emph{la siesta}, once more
  \item \emph{Say it:} a sentence of your own with \emph{la siesta}
  \item \emph{Recall:} say the Spanish for to fall asleep, then the Spanish for insomnia, then say \emph{la siesta} again
\end{itemize}

\begin{tcolorbox}[breakable,colback=teal!4,colframe=teal!35!black,title={Before you move on}]
What does la siesta mean? (\textquotedblleft{}The nap, the siesta\textquotedblright{}.) Say it once more.
\end{tcolorbox}
